\documentclass[11pt]{article}

\usepackage[margin=1in]{geometry}
\usepackage[T1]{fontenc}
\usepackage[utf8]{inputenc}
\usepackage{lmodern}
\usepackage{microtype}
\usepackage{amsmath, amssymb, mathtools}
\usepackage{bm}
\usepackage{booktabs}
\usepackage{enumitem}
\usepackage{hyperref}

\hypersetup{
  colorlinks=true,
  linkcolor=blue,
  urlcolor=blue
}

\setlist[itemize]{leftmargin=1.5em}
\setlength{\parskip}{0.6em}
\setlength{\parindent}{0pt}

\title{EFGPND Conditioning Note}
\author{}
\date{}

\begin{document}

\maketitle

This note updates the earlier $\sigma_n^2$ memo with the current picture:
\begin{itemize}
  \item the original noise-trace pathology was real, and the specialized
  feature-space formula fixed most of it;
  \item the remaining hard piece is now the shared mean solve
  \[
    (G + \sigma_n^2 I)\beta = D F^\ast y,
  \]
  not the old $\sigma_n^2$ trace block;
  \item this is primarily a large-$n$ conditioning problem;
  \item factoring out $1/\sigma_n^2$ by itself does \emph{not} improve the
  condition number of that mean system;
  \item the next algebraic opportunity for the remaining lengthscale problem is
  to differentiate the \emph{whole approximate objective} in feature space,
  rather than only rewriting the trace term.
\end{itemize}

\section*{1. Operators}

Write the approximate kernel in feature form as
\begin{equation}
\widetilde K = \Phi \Phi^\ast,
\qquad
\Phi = F D,
\end{equation}
and define
\begin{equation}
G = \Phi^\ast \Phi = D(F^\ast F)D.
\end{equation}

Then the data-space and feature-space systems are
\begin{equation}
K = \widetilde K + \sigma_n^2 I
= \Phi \Phi^\ast + \sigma_n^2 I,
\end{equation}
\begin{equation}
A_{\mathrm{mean}} = G + \sigma_n^2 I.
\end{equation}

The shared solve used to build
\[
\alpha = K^{-1}y
\]
is the feature-space system
\begin{equation}
A_{\mathrm{mean}} \beta = D F^\ast y,
\qquad
\alpha = \frac{1}{\sigma_n^2}\left(y - F D \beta\right).
\end{equation}

This same $\alpha$ enters the quadratic-form part of \emph{every}
hypergradient:
\begin{equation}
y^\top K^{-1} \left(\frac{\partial K}{\partial \theta}\right) K^{-1} y
=
\alpha^\top \left(\frac{\partial K}{\partial \theta}\right) \alpha.
\end{equation}

\section*{2. Why Large $n$ Makes This Bad}

The current difficulty is most apparent in large dense datasets such as the
PRISM raster because the top eigenvalues of $G = D(F^\ast F)D$ grow with the
sampling density. On a fixed domain with infill sampling, this means they grow
roughly with $n$, up to kernel-dependent constants and the scale set by
$\sigma_f^2$.

So the mean system has condition number
\begin{equation}
\kappa(A_{\mathrm{mean}})
\approx
\frac{\lambda_{\max}(G) + \sigma_n^2}{\sigma_n^2}.
\end{equation}

This becomes bad for three separate reasons:
\begin{itemize}
  \item large $n$ increases $\lambda_{\max}(G)$;
  \item large $\sigma_f^2$ increases the scale of $G$;
  \item small or moderate $\sigma_n^2$ leaves a weak floor on the bottom of the
  spectrum.
\end{itemize}

This explains why the problem is much more visible on PRISM than on smaller
datasets even when the number of features $M$ is only a few hundred.

\section*{3. Why Small Residual Does Not Mean Good Gradient}

CG controls the residual
\[
r = b - A \hat x,
\]
not the solution error directly. For a linear system $Ax = b$,
\begin{equation}
\frac{\|x - \hat x\|}{\|x\|}
\lesssim
\kappa(A)\,
\frac{\|r\|}{\|b\|}.
\end{equation}

So if $\kappa(A_{\mathrm{mean}})$ is on the order of $10^7$ to $10^8$, then a
relative residual tolerance such as $10^{-4}$ or even $10^{-5}$ is a weak
guarantee for the quality of $K^{-1}y$.

That is exactly what the fixed-state PRISM checks show: many CG iterations,
apparently small residuals, and still a poor approximation to the true
feature-space solve.

\section*{4. Why the $\sigma_n^2$ Fix Worked}

The earlier noise-variance problem was different. The old stochastic trace term
for $\sigma_n^2$ required repeated solves with harsh probe right-hand sides of
the form
\begin{equation}
A_{\mathrm{mean}}^{-1}(D F^\ast z).
\end{equation}

After the algebraic rewrite, the trace identity becomes
\begin{equation}
\operatorname{tr}(K^{-1})
=
\frac{n}{\sigma_n^2}
-
\frac{1}{\sigma_n^2}
\operatorname{tr}(A_{\mathrm{mean}}^{-1} G),
\end{equation}
and the corresponding Hutchinson estimator uses solves of the form
\begin{equation}
A_{\mathrm{mean}}^{-1}(Gv).
\end{equation}

The gain did \emph{not} come mainly from pulling out the factor
$1/\sigma_n^2$. The real gain came from changing the problem structurally:
\begin{itemize}
  \item an exact closed-form term $n/\sigma_n^2$ was peeled off;
  \item the remaining solves use $Gv$, so the right-hand side lives in
  $\mathrm{range}(G)$;
  \item weak or null directions are suppressed instead of being excited by raw
  white-noise probes.
\end{itemize}

That is why the $\sigma_n^2$ trace block became dramatically easier.

\section*{5. Does Factoring Out $1/\sigma_n^2$ Help the Mean Solve?}

Not by itself.

The code already has the scaled operator
\begin{equation}
A_{\mathrm{var}} = I + \frac{1}{\sigma_n^2} G,
\end{equation}
implemented as
\[
\texttt{create\_A\_var}.
\]

But
\begin{equation}
A_{\mathrm{mean}} = G + \sigma_n^2 I = \sigma_n^2 A_{\mathrm{var}}.
\end{equation}

Since multiplying a matrix by a positive scalar does not change its condition
number,
\begin{equation}
\kappa(A_{\mathrm{mean}})
=
\kappa(A_{\mathrm{var}}).
\end{equation}

So simply factoring out $1/\sigma_n^2$ does \emph{not} cure the mean-solve
conditioning problem. It only rescales the operator.

\section*{6. Current Empirical Picture on Late PRISM}

At a representative late PRISM state
\begin{equation}
\ell = 0.09256,
\qquad
\sigma_f^2 = 3.878,
\qquad
\sigma_n^2 = 0.05202,
\qquad
M = 625,
\end{equation}
I compared \texttt{compute\_gradients} against the exact feature-space solve for
the \emph{current approximate objective}, using the same Hutchinson probes.

The exact raw/log-parameter gradient was approximately
\begin{equation}
g_{\mathrm{exact}}
\approx
\begin{bmatrix}
1.40 \times 10^6 \\
-7.19 \times 10^4 \\
-6.36 \times 10^5
\end{bmatrix}.
\end{equation}

The observed relative raw-gradient errors were:

\begin{center}
\begin{tabular}{lccc}
\toprule
setting & mean CG iters & trace CG iters & relative raw-gradient error \\
\midrule
baseline, $\texttt{cg\_tol}=10^{-4}$ & 753 & 178 & $2.86 \times 10^2$ \\
preconditioned, $\texttt{cg\_tol}=10^{-4}$ & 166 & 50 & $3.21 \times 10^1$ \\
preconditioned + warm start, $\texttt{cg\_tol}=10^{-4}$ & 69 & 50 & $2.76 \times 10^1$ \\
preconditioned, $\texttt{cg\_tol}=10^{-5}$ & 977 & 194 & $3.22 \times 10^{-1}$ \\
preconditioned + warm start, $\texttt{cg\_tol}=10^{-5}$ & 685 & 194 & $2.21 \times 10^{-1}$ \\
\bottomrule
\end{tabular}
\end{center}

So:
\begin{itemize}
  \item preconditioning helps a lot;
  \item warm starting helps further;
  \item but late in training, even $\texttt{cg\_tol}=10^{-5}$ is still only
  moderately accurate, not exact.
\end{itemize}

\section*{7. The Next Algebraic Win: Special-Case $\sigma_f^2$}

The signal variance is also algebraically special. If
\begin{equation}
\widetilde K = \sigma_f^2 \widetilde K_0,
\end{equation}
then
\begin{equation}
\frac{\partial K}{\partial \sigma_f^2}
=
\frac{\widetilde K}{\sigma_f^2}
=
\frac{K - \sigma_n^2 I}{\sigma_f^2}.
\end{equation}

Therefore
\begin{equation}
\operatorname{tr}\!\left(K^{-1}\frac{\partial K}{\partial \sigma_f^2}\right)
=
\frac{1}{\sigma_f^2}
\operatorname{tr}\!\left(I - \sigma_n^2 K^{-1}\right)
=
\frac{n - \sigma_n^2 \operatorname{tr}(K^{-1})}{\sigma_f^2},
\end{equation}
and
\begin{equation}
\alpha^\top
\left(\frac{\partial K}{\partial \sigma_f^2}\right)
\alpha
=
\frac{1}{\sigma_f^2}
\left(
\alpha^\top K \alpha
- \sigma_n^2 \alpha^\top \alpha
\right)
=
\frac{y^\top \alpha - \sigma_n^2 \alpha^\top \alpha}{\sigma_f^2}.
\end{equation}

So the positive-parameter gradient for $\sigma_f^2$ is
\begin{equation}
\frac{\partial}{\partial \sigma_f^2}\log p(y \mid \theta)
=
\frac{1}{2\sigma_f^2}
\left(
n - \sigma_n^2 \operatorname{tr}(K^{-1})
- y^\top \alpha
+ \sigma_n^2 \alpha^\top \alpha
\right).
\end{equation}

In raw/log-parameter space, this becomes
\begin{equation}
g_{\log \sigma_f^2}
=
\sigma_f^2 \frac{\partial}{\partial \sigma_f^2}\log p(y \mid \theta)
=
\frac{1}{2}
\left(
n - \sigma_n^2 \operatorname{tr}(K^{-1})
- y^\top \alpha
+ \sigma_n^2 \alpha^\top \alpha
\right).
\end{equation}

Likewise, the raw/log-parameter noise gradient is
\begin{equation}
g_{\log \sigma_n^2}
=
\sigma_n^2 \frac{\partial}{\partial \sigma_n^2}\log p(y \mid \theta)
=
\frac{\sigma_n^2}{2}
\left(
\operatorname{tr}(K^{-1}) - \alpha^\top \alpha
\right).
\end{equation}

This means that once $\alpha$ and $\operatorname{tr}(K^{-1})$ are available,
both $\sigma_f^2$ and $\sigma_n^2$ can be handled \emph{without} their own
stochastic derivative blocks. Only the lengthscale still needs the hard
derivative-specific trace machinery.

\section*{8. Methodological Implication}

The main lesson is not merely ``use more CG iterations.'' The better
methodology is:
\begin{itemize}
  \item exploit algebraic structure for special hyperparameters whenever
  possible;
  \item avoid stochastic derivative solves for $\sigma_n^2$ and $\sigma_f^2$;
  \item recognize that the remaining hard core is the shared mean solve and the
  lengthscale derivative, both of which are limited by conditioning rather than
  by missing algebra.
\end{itemize}

So the most promising path forward is:
\begin{enumerate}[leftmargin=1.5em]
  \item keep the specialized $\sigma_n^2$ treatment;
  \item special-case $\sigma_f^2$ as above;
  \item then focus numerical work on the mean solve
  $(G + \sigma_n^2 I)\beta = D F^\ast y$, where large-$n$ conditioning is now
  the real bottleneck.
\end{enumerate}

\section*{9. Lengthscale: What Failed}

After special-casing $\sigma_n^2$ and $\sigma_f^2$, the remaining
hyperparameter-specific difficult block is the lengthscale trace term
\[
\frac{1}{2}\operatorname{tr}\!\left(K^{-1}\partial_\ell K\right).
\]

The first idea was the direct analogue of the successful $\sigma_n^2$ rewrite.
Let
\begin{equation}
C = F^\ast F,
\qquad
D = \operatorname{diag}(w),
\qquad
A = D C D + \sigma_n^2 I,
\end{equation}
and write
\begin{equation}
\partial_\ell \widetilde K = F S_\ell F^\ast,
\qquad
S_\ell = \operatorname{diag}(d_\ell),
\qquad
d_\ell = \frac{\partial (w^2)}{\partial \ell}.
\end{equation}

Using the Woodbury form
\begin{equation}
K^{-1}
=
\sigma_n^{-2}\left(I - \Phi A^{-1}\Phi^\ast\right),
\qquad
\Phi = F D,
\end{equation}
one gets the exact identity
\begin{equation}
\operatorname{tr}\!\left(K^{-1}\partial_\ell K\right)
=
\frac{1}{\sigma_n^2}\operatorname{tr}(C S_\ell)
-
\frac{1}{\sigma_n^2}
\operatorname{tr}\!\left(A^{-1} D C S_\ell C D\right).
\end{equation}

This identity is correct, but the resulting estimator is not numerically useful
in the low-probe regime of interest.

\subsection*{Why It Failed in Practice}

The problem is that the full lengthscale derivative is \emph{sign-changing}.
So the formula above is an exact constant minus a noisy stochastic correction
for an indefinite operator. At $J=1$, this creates very large cancellation
variance.

On a small exact check with real data ($N=192$, $M=625$), the true lengthscale
trace was about
\[
\operatorname{tr}\!\left(K^{-1}\partial_\ell K\right) \approx -3.576\times 10^3.
\]
Over 1000 seeds with $J=1$:
\begin{itemize}
  \item the old estimator had mean/std approximately
  \[
    -3.575\times 10^3 \;\; / \;\; 7.10\times 10^2;
  \]
  \item the full feature-space rewrite had mean/std approximately
  \[
    -3.055\times 10^3 \;\; / \;\; 8.82\times 10^4.
  \]
\end{itemize}

So the direct rewrite is algebraically right, but statistically terrible at
$J=1$.

The same pattern showed up on the frozen late-PRISM state used above
($\ell=0.09256$, $\sigma_f^2=3.878$, $\sigma_n^2=0.05202$, $M=625$). Over 20
seeds with $J=1$:
\begin{itemize}
  \item the old lengthscale trace estimator had sample mean/std approximately
  \[
    -7.38\times 10^5 \;\; / \;\; 1.17\times 10^6;
  \]
  \item the full rewrite had sample mean/std approximately
  \[
    -1.13\times 10^9 \;\; / \;\; 1.01\times 10^{10}.
  \]
\end{itemize}

That is why it destabilized training even though the underlying identity was
correct.

\section*{10. A Better Route: Exact Scale Term + PSD Residual}

The more promising route is \emph{not} to rewrite the whole lengthscale
derivative at once. Instead, split it into:
\begin{itemize}
  \item an exact amplitude-like piece proportional to $\widetilde K$;
  \item a residual piece that is positive semidefinite in feature space.
\end{itemize}

Suppose the lengthscale derivative of the spectral weights can be written as
\begin{equation}
\frac{\partial (w^2)}{\partial \ell}
=
c_\ell \, w^2 - q_\ell,
\qquad
q_\ell(\xi) \ge 0
\;\;\text{for all }\xi.
\end{equation}
Equivalently,
\begin{equation}
\partial_\ell \widetilde K
=
c_\ell \widetilde K - F Q_\ell F^\ast,
\qquad
Q_\ell = \operatorname{diag}(q_\ell).
\end{equation}

Then the trace term becomes
\begin{align}
\operatorname{tr}\!\left(K^{-1}\partial_\ell K\right)
&=
c_\ell \operatorname{tr}(K^{-1}\widetilde K)
-
\operatorname{tr}\!\left(K^{-1} F Q_\ell F^\ast\right) \\
&=
c_\ell \left(n - \sigma_n^2 \operatorname{tr}(K^{-1})\right)
-
\operatorname{tr}\!\left(K^{-1} F Q_\ell F^\ast\right).
\end{align}

The first term is exact once $\operatorname{tr}(K^{-1})$ is known from the
stable $\sigma_n^2$ machinery.

For the residual term, the same Woodbury manipulation gives
\begin{equation}
\operatorname{tr}\!\left(K^{-1} F Q_\ell F^\ast\right)
=
\frac{1}{\sigma_n^2}\operatorname{tr}(C Q_\ell)
-
\frac{1}{\sigma_n^2}
\operatorname{tr}\!\left(A^{-1} D C Q_\ell C D\right).
\end{equation}

This has the same basic structure as the successful noise rewrite, but now the
operator being estimated is associated with the \emph{positive} residual
$Q_\ell$, not the full sign-changing derivative.

\subsection*{Why This Is More Plausible}

This route avoids the thing that killed the full rewrite: it does not ask one
Hutchinson estimator to represent the entire indefinite lengthscale derivative
as a single constant-minus-noisy correction.

Instead:
\begin{itemize}
  \item the exact amplitude-like piece is peeled off deterministically;
  \item the stochastic part is only the residual PSD operator
  \[
    F Q_\ell F^\ast;
  \]
  \item the feature-space solve still uses the same mean-system matrix
  \[
    A = D C D + \sigma_n^2 I.
  \]
\end{itemize}

That is much closer in spirit to the $\sigma_n^2$ rewrite, which was stable
because it estimated a filtered PSD object rather than a sign-changing one.

\section*{11. SE and Mat\'ern Both Fit This Pattern}

\subsection*{Squared Exponential}

For the squared exponential kernel,
\begin{equation}
S(\xi)
=
\sigma_f^2
\left((2\pi)\ell^2\right)^{d/2}
\exp\!\left(-2\pi^2 \ell^2 \|\xi\|^2\right),
\end{equation}
so
\begin{equation}
\frac{\partial S}{\partial \ell}
=
S(\xi)
\left(
\frac{d}{\ell}
- 4\pi^2 \ell \|\xi\|^2
\right).
\end{equation}

Therefore
\begin{equation}
\frac{\partial (w^2)}{\partial \ell}
=
\frac{d}{\ell} w^2
-
\underbrace{4\pi^2 \ell \|\xi\|^2 w^2}_{q_\ell(\xi) \ge 0}.
\end{equation}

So the SE kernel fits the split exactly with
\begin{equation}
c_\ell = \frac{d}{\ell},
\qquad
Q_\ell = \operatorname{diag}\!\left(4\pi^2 \ell \|\xi\|^2 w^2\right).
\end{equation}

\subsection*{Mat\'ern}

For the Mat\'ern family in the parameterization used here,
\begin{equation}
S(\xi)
=
\sigma_f^2 c_{\nu,d}\,
\ell^d
\left(2\nu + 4\pi^2 \ell^2 \|\xi\|^2\right)^{-(\nu + d/2)},
\end{equation}
for a positive constant $c_{\nu,d}$ independent of $\ell$ and $\xi$.
Differentiating gives
\begin{equation}
\frac{\partial S}{\partial \ell}
=
S(\xi)
\left(
\frac{d}{\ell}
-
\frac{8\pi^2(\nu + d/2)\,\ell \|\xi\|^2}
{2\nu + 4\pi^2 \ell^2 \|\xi\|^2}
\right).
\end{equation}

So again
\begin{equation}
\frac{\partial (w^2)}{\partial \ell}
=
\frac{d}{\ell} w^2
-
\underbrace{
\frac{8\pi^2(\nu + d/2)\,\ell \|\xi\|^2}
{2\nu + 4\pi^2 \ell^2 \|\xi\|^2}
w^2
}_{q_\ell(\xi) \ge 0}.
\end{equation}

Therefore the same exact-plus-residual strategy is not just SE-specific; it
extends naturally to Mat\'ern as well.

\section*{12. What Computation Would Actually Change}

The key term to change is still only
\[
\frac{1}{2}\operatorname{tr}\!\left(K^{-1}\partial_\ell K\right).
\]

Under the proposed split, this becomes:
\begin{itemize}
  \item an exact scale term
  \[
    \frac{c_\ell}{2}\left(n - \sigma_n^2 \operatorname{tr}(K^{-1})\right),
  \]
  \item plus a residual trace term built from
  \[
    H_\ell^{\mathrm{res}} = D C Q_\ell C D.
  \]
\end{itemize}

The stochastic residual is still a Hutchinson trace estimate in feature space:
\begin{equation}
\operatorname{tr}\!\left(A^{-1} H_\ell^{\mathrm{res}}\right)
\approx
\frac{1}{J}\sum_{j=1}^J v_j^\ast A^{-1} H_\ell^{\mathrm{res}} v_j.
\end{equation}

So there is still:
\begin{itemize}
  \item a stochastic trace term;
  \item a CG solve with the same operator
  \[
    A = D C D + \sigma_n^2 I;
  \]
  \item one solve per probe.
\end{itemize}

What changes is the right-hand side construction. The residual right-hand side
is
\begin{equation}
b_j^{\mathrm{res}}
=
H_\ell^{\mathrm{res}} v_j
=
D C Q_\ell C D v_j,
\end{equation}
which uses only Toeplitz applications of $C = F^\ast F$ and diagonal
multiplications.

\section*{13. Complexity and Expected Benefit}

This residual rewrite would \emph{not} fix the shared mean solve
\[
(G + \sigma_n^2 I)\beta = D F^\ast y,
\]
which is still the main common bottleneck.

But it would make the lengthscale-specific trace work look like
\begin{equation}
O\!\left(J M \log m\right)
+
O\!\left(k_{\mathrm{cg}} J M \log m\right),
\end{equation}
up to moderate constants, while keeping the stochastic part tied to a PSD
residual operator rather than the full sign-changing derivative.

So the best current picture is:
\begin{itemize}
  \item the full direct lengthscale rewrite is mathematically correct but too
  noisy at low $J$;
  \item the natural next attempt is a \emph{partial} rewrite:
  exact scale term plus PSD residual;
  \item that route generalizes cleanly to SE and Mat\'ern, but in its naive form
  it is still not good enough at $J=1$.
\end{itemize}

\subsection*{Current Status of the PSD-Residual Split}

I tested this split directly against exact ground truth for the approximate
EFGP objective, both on small dense real-data subsets and on the full PRISM
state with exact $M\times M$ systems.

The result is disappointing: the naive split is still dramatically noisier than
the old estimator in the $J=1$ regime.

On the small SE check ($N=192$, $M=625$), the exact lengthscale trace was about
\[
-3.576\times 10^3.
\]
Over 1000 seeds with $J=1$:
\begin{itemize}
  \item the old estimator had mean/std approximately
  \[
    -3.594\times 10^3 \;\; / \;\; 7.00\times 10^2;
  \]
  \item the PSD-residual split had mean/std approximately
  \[
    -1.506\times 10^3 \;\; / \;\; 1.24\times 10^5.
  \]
\end{itemize}

On the small Mat\'ern check ($N=192$, $M=3481$):
\begin{itemize}
  \item the old estimator had mean/std approximately
  \[
    -1.973\times 10^3 \;\; / \;\; 2.06\times 10^2;
  \]
  \item the PSD-residual split had mean/std approximately
  \[
    -2.464\times 10^3 \;\; / \;\; 4.33\times 10^4.
  \]
\end{itemize}

On the frozen full-PRISM state, the exact lengthscale trace was about
\[
-3.628\times 10^4.
\]
Over 20 seeds with $J=1$:
\begin{itemize}
  \item the old estimator had mean/std approximately
  \[
    -3.606\times 10^4 \;\; / \;\; 7.65\times 10^3;
  \]
  \item the PSD-residual split had mean/std approximately
  \[
    -4.673\times 10^8 \;\; / \;\; 1.95\times 10^9.
  \]
\end{itemize}

So the practical conclusion is now:
\begin{itemize}
  \item the full rewrite fails;
  \item the exact-scale-plus-PSD-residual split is conceptually cleaner, but
  still fails badly as a low-probe drop-in estimator;
  \item positive semidefiniteness of the residual alone is not enough to make
  the resulting Hutchinson estimator usable at $J=1$.
\end{itemize}

That means the next improvement for lengthscale probably has to come from a
different direction: better probe design, correlated probes, control variates,
or a different optimization strategy for the lengthscale parameter itself,
rather than another straightforward constant-plus-feature-space-trace rewrite.

\section*{14. Whole-Objective Feature-Space Derivative}

The next algebraic candidate is to stop splitting the lengthscale gradient into
``trace term'' and ``quadratic term'' in data space, and instead differentiate
the entire approximate objective directly in feature space.

Write
\begin{equation}
K = \sigma_n^2 I + \Phi \Phi^\ast,
\qquad
\Phi = F D,
\qquad
G = \Phi^\ast \Phi = D C D,
\qquad
C = F^\ast F.
\end{equation}
Define also
\begin{equation}
A = G + \sigma_n^2 I,
\qquad
b = D F^\ast y.
\end{equation}

Then the approximate negative log marginal likelihood can be written exactly as
\begin{equation}
\mathcal L(\theta)
=
\frac{1}{2}\left[
(n-M)\log \sigma_n^2
\;+\;
\log\det A
\;+\;
\frac{1}{\sigma_n^2}
\left(
y^\top y - b^\ast A^{-1} b
\right)
\right]
 + \mathrm{const}.
\end{equation}

Now consider any hyperparameter $\theta$ that changes only the spectral weights
$D$ and not $\sigma_n^2$. Then
\begin{equation}
\partial_\theta \mathcal L
=
\frac{1}{2}\left[
\operatorname{tr}(A^{-1} A_\theta)
-
\frac{1}{\sigma_n^2}
\left(
2\,\Re(b_\theta^\ast \beta)
-
\beta^\ast A_\theta \beta
\right)
\right],
\qquad
\beta = A^{-1} b.
\end{equation}

Here
\begin{equation}
A_\theta = D_\theta C D + D C D_\theta,
\qquad
b_\theta = D_\theta F^\ast y.
\end{equation}

It is better numerically to avoid differentiating $w$ directly. Since
$w^2 = s$, define
\begin{equation}
E_\theta
=
\operatorname{diag}\!\left(
\frac{1}{2}\frac{\partial_\theta s}{s}
\right),
\end{equation}
so that $D_\theta = E_\theta D$ and therefore
\begin{equation}
A_\theta = E_\theta G + G E_\theta,
\qquad
b_\theta = E_\theta b.
\end{equation}

\subsection*{What Changes Computationally}

The old lengthscale path uses
\begin{equation}
\frac{1}{2}
\left(
\operatorname{tr}(K^{-1}\partial_\ell K)
-
\alpha^\ast (\partial_\ell K)\alpha
\right),
\end{equation}
with a data-space Hutchinson estimator for the trace term and the quadratic
term expressed through $\alpha = K^{-1} y$.

The proposed feature-space objective derivative instead uses
\begin{equation}
\frac{1}{2}\left[
\operatorname{tr}(A^{-1} A_\ell)
-
\frac{1}{\sigma_n^2}
\left(
2\,\Re(b_\ell^\ast \beta)
-
\beta^\ast A_\ell \beta
\right)
\right].
\end{equation}

Only the trace part remains stochastic:
\begin{equation}
\operatorname{tr}(A^{-1} A_\ell)
\approx
\frac{1}{J}\sum_{j=1}^J v_j^\ast A^{-1} A_\ell v_j,
\end{equation}
with feature-space Rademacher probes $v_j \in \{\pm 1\}^M$.

The right-hand side is now
\begin{equation}
A_\ell v
=
E_\ell G v + G(E_\ell v),
\end{equation}
which requires only Toeplitz applications of $G = D C D$ and diagonal
multiplications.

\subsection*{Why This Might Be Better}

Relative to the older lengthscale trace estimator, this route has three
potential advantages:
\begin{itemize}
  \item it keeps the entire derivative in feature space;
  \item it avoids the large deterministic-constant minus noisy-correction
  structure that killed the previous trace-only rewrites;
  \item it replaces the data-space derivative probe path with Toeplitz-based
  applications of $A_\ell$.
\end{itemize}

But there is also a real risk: $A_\ell$ is still generally sign-indefinite, so
the Hutchinson trace estimator may still have high variance at very small $J$.
That has to be checked empirically, not assumed.

\subsection*{Empirical Outcome}

I tested this whole-objective feature-space estimator against exact ground truth
for the approximate EFGP objective in the same style as the earlier
lengthscale experiments:
\begin{itemize}
  \item small real-data subsets with explicit dense $K$ available;
  \item the frozen late-PRISM state using exact formed $M\times M$ systems.
\end{itemize}

The algebra checks out. On the small dense checks, the exact feature-space
gradient matched the exact dense-$K$ gradient to about $2\times 10^{-3}$ in
absolute error.

But as an estimator, the picture is mixed:
\begin{itemize}
  \item it is \emph{much} better than the previous failed feature-space
  rewrites;
  \item it is still materially noisier than the old estimator at $J=1$;
  \item its feature-space trace solve is harder for CG than the old
  lengthscale trace solve.
\end{itemize}

On the small SE check ($N=192$, $M=625$), with $J=1$ over 1000 seeds:
\begin{itemize}
  \item the old estimator had gradient mean/std approximately
  \[
    3.268\times 10^2 \;\; / \;\; 3.500\times 10^2;
  \]
  \item the whole-objective feature-space estimator had gradient mean/std
  approximately
  \[
    1.834\times 10^2 \;\; / \;\; 4.611\times 10^3.
  \]
\end{itemize}

On the small Mat\'ern check ($N=192$, $M=3481$):
\begin{itemize}
  \item the old estimator had gradient mean/std approximately
  \[
    -1.987\times 10^2 \;\; / \;\; 1.031\times 10^2;
  \]
  \item the whole-objective feature-space estimator had gradient mean/std
  approximately
  \[
    -1.444\times 10^2 \;\; / \;\; 1.678\times 10^3.
  \]
\end{itemize}

So on the small checks it is clearly worse than the current estimator.

On the frozen late-PRISM state ($N \approx 1.21\times 10^7$, $M=625$), with
$J=1$ over 20 seeds:
\begin{itemize}
  \item the old estimator had gradient mean/std approximately
  \[
    1.507167\times 10^7 \;\; / \;\; 3.827\times 10^3;
  \]
  \item the whole-objective feature-space estimator had gradient mean/std
  approximately
  \[
    1.512696\times 10^7 \;\; / \;\; 6.854\times 10^5.
  \]
\end{itemize}

This is still far noisier than the old estimator, but it is no longer
catastrophically bad. The relative standard deviation is about $4.5\%$ on this
PRISM state, versus about $2.5\times 10^{-4}$ for the old estimator.

The speed story goes the other way. On the same frozen PRISM state with the
actual fast operators and Jacobi preconditioning:
\begin{itemize}
  \item old lengthscale path: about $1.23$s, with about $172$ CG iterations;
  \item whole-objective feature-space path: about $0.17$s, but about $1250$ CG
  iterations, i.e.\ it hits the default $2M$ cap.
\end{itemize}

So the new estimator removes almost all of the $N$-scale work, but it replaces
it with a substantially harder feature-space Krylov problem.

The practical conclusion is:
\begin{itemize}
  \item this is \emph{not} a clean drop-in replacement for the current
  lengthscale estimator;
  \item it is algebraically sound and dramatically cheaper in wall clock on
  very large datasets;
  \item but at low probe count it is still noisier than the old estimator, and
  its trace solve is harder for CG;
  \item if this direction is pursued further, it probably needs additional
  variance reduction or a stronger/recycled solver, not just the bare algebraic
  rewrite.
 \end{itemize}

\end{document}
